\documentclass[11pt]{article}
\usepackage[letterpaper,margin=0.82in]{geometry}
\usepackage[T1]{fontenc}
\usepackage{array}
\usepackage{longtable}
\usepackage{hyperref}

\title{\textbf{TensorGuard: Static Refinement Verification for PyTorch Tensor Programs}}
\author{TensorGuard Authors}
\date{\today}

\newcommand{\tg}{TensorGuard}
\newcommand{\shape}{\ensuremath{\mathsf{Shape}}}
\newcommand{\device}{\ensuremath{\mathsf{Device}}}
\newcommand{\phase}{\ensuremath{\mathsf{Phase}}}
\newcommand{\dtype}{\ensuremath{\mathsf{DType}}}
\newcommand{\grad}{\ensuremath{\mathsf{Grad}}}
\newcommand{\unknown}{\ensuremath{\mathsf{UNKNOWN}}}
\newcommand{\safe}{\ensuremath{\mathsf{SAFE}}}
\newcommand{\unsafe}{\ensuremath{\mathsf{UNSAFE}}}
\setlength{\parskip}{0.45em}
\setlength{\parindent}{0pt}

\begin{document}
\maketitle

\begin{abstract}
\tg{} is a static verifier for PyTorch models that turns tensor-program
well-formedness into refinement constraints.  It checks shape, device, dtype,
train/eval phase, stride, permutation, and gradient-flow facts before training
or export begins, while preserving an explicit three-valued soundness contract:
when the verifier cannot prove a claim inside its supported fragment it reports
\unknown{} rather than silently passing.  The implementation combines a source
frontend, a \texttt{torch.fx}/export-oriented graph frontend, a reduced product
of tensor abstract domains, Z3-backed safety checks, CEGAR-style predicate
discovery, differential conformance tests against live PyTorch, Lean proof
artifacts for core rules, and reproducible empirical harnesses.

The repository is not merely a prototype pass over toy models.  It contains
operator transfer functions for everyday PyTorch programs, frontends for real
\texttt{nn.Module} code, integration hooks for CLIs, GitHub Actions, SARIF,
watch mode, Jupyter, decorators, framework adapters, deployment gates, and
artifact scripts that regenerate precision/recall, latency, ablation, coverage,
conformance, fuzzing, mutation, and benchmark evidence.  This paper describes
the system as shipped in the repository, with emphasis on the engineering and
evaluation choices that make \tg{} a credible research artifact and a useful
developer tool.
\end{abstract}

\clearpage

\section{Problem and Claim}

Most deep-learning correctness tools see tensor bugs too late.  A model can
import successfully, type-check as Python, instantiate its modules, and even
reach a training loop before a wrong channel count, a bad attention projection,
an accidental CPU buffer, a mismatched dtype, or a train/eval phase dependency
raises a runtime exception.  Worse, some failures are silent: a detached tensor
can sever gradient flow, a broadcast can mask an intended batch axis, or an
export pipeline can accept a graph whose serving contract no longer matches the
model author's intent.

\tg{} attacks these failures as static contract violations.  A tensor value is
modeled with refinement information such as
\[
  \{v : \mathrm{Tensor} \mid \shape(v)=(B,C,H,W)
       \wedge \device(v)=\mathrm{cpu}
       \wedge \dtype(v)=\mathrm{float32}\}.
\]
The verifier propagates these facts through constructors, module calls,
functional operators, reshapes, broadcasts, indexing, attention blocks,
recurrent modules, normalization layers, torchvision v2 image transforms,
sparse layouts and sparse-dense kernels, distributed wrappers, and deployment
gates.  At each operation it asks whether the preconditions of the transfer
function are entailed by the current abstract state.  Concrete contradictions
become high-confidence bugs.  Unknown or out-of-fragment regions become
abstentions with reasons, not fake successes.

The central claim of the repository is practical: a tensor verifier can be
soundly conservative and still catch real bugs early across modern PyTorch
workflows.  The codebase supports that claim with four pillars.

\begin{enumerate}
\item A precise public contract that separates \safe{}, \unsafe{}, and
\unknown{} outcomes, with a strict ``sound'' mode for CI.
\item A broad operator and frontend implementation that handles real modules,
not only a calculus in a paper.
\item A reproducible evaluation suite with real-bug corpora, differential
testing, mutation tests, benchmark scripts, ablations, dashboards, and baseline
comparisons.
\item A formalization and conformance story: Lean artifacts for core transfer
rules, Z3 encodings for symbolic constraints, and live PyTorch oracles for
operator behavior.
\end{enumerate}

This is a tool paper: the interesting contribution is the system boundary,
not a single transfer function.  The repository demonstrates how a static
analysis tool can be packaged so that reviewers, users, and future contributors
can inspect the guarantee, reproduce the headline evidence, and extend the
coverage without weakening the trust story.

\section{Architecture}

\tg{} is organized as a pipeline from Python source or captured graphs to a
multi-domain safety decision.  The source path parses \texttt{nn.Module}
classes, extracts constructor parameters and layer definitions, walks
\texttt{forward}, and emits a computation graph.  The graph path consumes
\texttt{torch.fx}, \texttt{torch.export}, and integration-specific traces where
available.  Both paths feed a common verifier built around symbolic tensor
states.

Each state maps tensor names to abstract facts.  Shape facts are tuples of
concrete integers or symbolic names; device facts range over CPU and CUDA
devices; phase facts track train/eval context; dtype facts model promotion and
casts; gradient facts track whether a value can carry gradient to downstream
parameters.  Several later steps add stride and permutation information so
viewability and layout-sensitive operators can be checked without collapsing
to opaque shapes.  These domains form a reduced product: a violation in one
domain can be explained with facts from the others, while disabled domains are
gated out of the solver rather than filtered only at the end.

The pipeline has four visible outputs.  First, it returns structured bugs with
locations, severities, confidence values, and optional counterexamples.  Second,
it returns a verdict: \safe{}, \unsafe{}, or \unknown{}.  Third, it can emit
developer-facing diagnostics, including source snippets, inferred-vs-expected
shapes, explanations, and autofix suggestions for mechanical cases.  Fourth,
it can emit machine-readable reporters: JSON, JUnit XML, SARIF 2.1.0, GitHub
annotations, coverage artifacts, dashboards, and certificates.

\subsection{Frontends}

The source frontend is deliberately useful on unannotated Python.  It recognizes
\texttt{nn.Linear}, convolution families, pooling, normalization, attention,
activation layers, containers, helper methods, subclassing patterns, and
functional calls such as \texttt{F.linear}, \texttt{F.pad},
\texttt{F.interpolate}, and \texttt{F.scaled\_dot\_product\_attention}.  It records scalar constructor
attributes such as head counts and hidden sizes, so common model code like
\texttt{x.view(B, T, self.num\_heads, D)} can be checked against the same
constants used by the layers.

The graph frontends target execution-oriented workflows.  \texttt{torch.fx}
support traces real modules and preserves enough source mapping for diagnostics.
\texttt{torch.export} and Dynamo-oriented paths let \tg{} run as a pre-pass
for export and compile pipelines.  The repository also includes integration
points for guarded \texttt{torch.compile}, AOTInductor packaging, ONNX export,
GGUF/llama.cpp-oriented checks, framework hooks, and deployment request/response
schemas.  These paths are important because many production tensor bugs appear
only when a model is compiled, exported, quantized, sharded, or served.

\subsection{Verifier}

The verifier performs stepwise symbolic propagation.  For each computation
step it updates the tensor state, emits safety conditions, and asks Z3 whether
a violation is forced.  Concrete mismatches, such as a \texttt{Linear} layer
expecting 512 features but receiving 256, are immediate.  Symbolic mismatches
are handled conservatively: if the verifier can prove incompatibility under
the known assumptions it reports a bug; if the result depends on an unconstrained
symbol, it abstains or keeps the dimension symbolic.

The same engine also supports CEGAR.  When the current abstraction is too weak,
\tg{} proposes additional shape predicates, reruns the proof attempt, and
records the discovered predicates as metadata.  If the predicates jointly imply
an impossible input contract, the conflict is promoted to a real
\texttt{cegar\_refined\_contract} bug rather than hidden as an internal note.
The repository includes convergence measurements and Lean-checked statements
for the predicate-bound argument.

\subsection{Soundness Modes}

The user-facing contract is intentionally explicit.  In \texttt{sound} mode,
\tg{} reports \safe{} only when the program remains inside the verifiable
fragment and all relevant transfer functions are at an acceptable confidence
level.  Unsupported layers, opaque data-dependent control flow, or heuristic
operators make the verdict \unknown{} unless a concrete bug is found.  In
\texttt{balanced} mode the tool remains useful for day-to-day development by
allowing more best-effort reasoning, while still surfacing abstentions.  In
\texttt{heuristic} mode users can opt into permissive exploratory checks.
The mode changes the verdict boundary; it does not hide discovered bugs.

\clearpage

\section{Tensor Semantics and Transfer Functions}

A transfer function is the contract for one operator family.  It consumes input
abstract facts and produces output facts or a violation.  \tg{} tags each
operator as complete, sound, or heuristic.  Complete rules are exact
shape-preserving families such as pointwise activations.  Sound rules enforce
well-defined structural contracts, even if the implementation abstains on hard
symbolic cases.  Heuristic rules are used when output rank or size depends on
runtime values, or when the rule intentionally approximates a large behavior
surface.

The repository includes transfer functions and tests for a wide range of
PyTorch constructs.  The core structural set covers matrix multiplication,
batched matrix multiplication, linear layers, convolution and transpose
convolution families, pooling, flatten, reshape/view with product reasoning,
broadcasting, reductions, cat/stack-family aliases, PyTorch-faithful padding,
interpolation/upsampling, transpose/permute, squeeze/unsqueeze,
\texttt{movedim}/\texttt{swapaxes}, \texttt{roll}/\texttt{rot90}/\texttt{flip},
\texttt{repeat\_interleave}/\texttt{tile}, tensor broadcasting helpers,
fold/unfold image-column transforms,
embedding, dtype casts, device moves, and stochastic tensor factories whose
shapes are seed-independent.  Extended coverage includes einsum parsing,
advanced indexing, gather/scatter, index select, masked select,
\texttt{take\_along\_dim}, \texttt{argsort}, \texttt{sort}/\texttt{topk}/
\texttt{kthvalue}, arg reductions, exact \texttt{nn/F.embedding\_bag}
contracts for offsets, per-sample weights, pooling modes, and TorchRec-style
jagged pooled embeddings, sparse COO/CSR/CSC/BSR/BSC layout invariants,
sparse-dense \texttt{mm}/\texttt{addmm}, CSR \texttt{sampled\_addmm},
sparse softmax/coalesce, layout conversions, FFT/complex view contracts, exact \texttt{torch.linalg} contracts for
\texttt{solve}, inverse, Cholesky, SVD, eigendecomposition, and QR, named tensor
alignment, vmap/functorch batch dimensions, \texttt{torch.func} autodiff
wrappers (\texttt{grad}, \texttt{jacrev}, \texttt{jacfwd}, \texttt{jvp}, and
\texttt{vjp}), distribution batch/event/log-prob
shapes, grid sampling, affine grids, source-level
\texttt{torchvision.transforms.v2} tensor transforms
(\texttt{Resize}, crops, \texttt{Pad}, \texttt{Normalize}, flips, and
PIL-only abstention), multi-head attention, scaled dot-product
attention, training loss contracts for \texttt{CrossEntropyLoss},
\texttt{NLLLoss}, \texttt{MSELoss}, \texttt{BCEWithLogitsLoss}, and
\texttt{KLDivLoss}, exact RNN/LSTM/GRU output and returned-state contracts,
LoRA adapter compatibility, training-loop checks, quantization/export safety,
mixed-precision/autocast dtype-divergence gates, distributed sharding, tensor
parallel attention variants, and deployment schemas.

Recent stack-family work follows the same standard.  \texttt{torch.stack},
\texttt{hstack}, \texttt{vstack}, \texttt{dstack},
\texttt{column\_stack}, and \texttt{row\_stack} are modeled with PyTorch's
singleton and rank-promotion corner cases, then checked by tests that run the
real operators and compare both output shapes and static error parity.

Axis-structural operators are now treated with the same precision rather than
hidden behind a generic shape-preserving fallback.  \texttt{squeeze} and
\texttt{unsqueeze} use PyTorch's exact rank rules, including tuple-dim
\texttt{squeeze} and the legal \([-\mathrm{rank}-1,\mathrm{rank}]\)
\texttt{unsqueeze} interval.  \texttt{movedim}/\texttt{moveaxis} use the ATen
axis-order algorithm, so multi-axis placements preserve symbolic dimension
identity instead of merely preserving rank.  \texttt{swapaxes}/\texttt{swapdims}
share the transpose transfer.  \texttt{roll}, \texttt{rot90}, and \texttt{flip}
preserve or permute shapes while still refuting statically invalid axis
contracts such as duplicate flip axes, mismatched roll arity, or invalid rotation
axes.  The regression suite compares source and FX extraction against live
PyTorch, exhaustively checks rank-4 \texttt{movedim} placements, and proves that
symbolic dimensions do not create false structural bugs.

Repeat and broadcast helpers are modeled as first-class structural transfers,
not as generic repeats or activations.  \texttt{repeat} and \texttt{tile} use
their distinct eager-PyTorch rank rules: \texttt{repeat} requires enough factors
for the input rank, while \texttt{tile} left-pads factors before multiplying
axes.  \texttt{repeat\_interleave} preserves the flattening behavior of
\texttt{dim=None}, statically checks \texttt{output\_size} when the repeated
extent is provable, and otherwise keeps a fresh symbolic extent instead of
claiming a false proof for data-dependent repeat vectors.
\texttt{torch.broadcast\_tensors} now gives every returned tensor the exact joint
broadcast shape, while source-level \texttt{torch.broadcast\_shapes} is carried
as a shape value that can feed \texttt{expand}.  Focused source and FX tests
compare these contracts against live PyTorch successes and exceptions.

Sparse tensors are checked beyond construction.  \tg{} verifies sparse-dense
\texttt{torch.sparse.mm} and \texttt{addmm} for COO/CSR/CSC matrices, exact
one-way \texttt{addmm} input broadcasting, CSR \texttt{sampled\_addmm} including
batched dense operands, COO-only sparse softmax and coalescing,
\texttt{to\_dense}, and \texttt{to\_sparse\_\{coo,csr,csc,bsr,bsc\}}
conversions with blocked-layout divisibility.  The parity tests compare every
successful output to live PyTorch via dense shape equality and keep unsupported
BSR/BSC sparse matrix kernels as explicit non-passes rather than build-dependent
silent approvals.

\subsection{Exact Index-Value Operators}

One recent implementation round strengthens a family that returns
indices, values, or both: \texttt{torch.take\_along\_dim}, \texttt{argsort},
\texttt{sort}, \texttt{topk}, \texttt{kthvalue}, \texttt{argmax}, and
\texttt{argmin}.  These operators are small but security-critical for a static
shape tool because the index tensors often feed a later gather.  Losing the fact
that an unpacked \texttt{indices} output is \texttt{int64}, or treating
\texttt{topk} as shape-preserving, can hide a real downstream error.

\tg{} now models the PyTorch contracts directly.  \texttt{take\_along\_dim}
returns the index shape when a dimension is supplied, checks same-rank and
non-dimension broadcastability, rejects non-\texttt{int64} indices, and flattens
to the index numel when \texttt{dim=None}.  \texttt{argsort}, \texttt{argmax},
and \texttt{argmin} produce \texttt{int64} outputs with the same or reduced
shape according to \texttt{keepdim}.  \texttt{sort}, \texttt{topk}, and
\texttt{kthvalue} propagate value tensors with the input dtype and index tensors
with \texttt{int64}, including FX \texttt{getitem}, namedtuple
\texttt{.values}/\texttt{.indices}, tensor-method calls, functional calls, and
source-level tuple unpacking.  Static \texttt{k} errors are refuted before
runtime.  The tests run real PyTorch tensors and compare both successful shapes
and runtime exceptions.

The same principle now covers the gather/scatter family.  \texttt{gather},
\texttt{scatter}/\texttt{scatter\_add}, and \texttt{index\_select} distinguish
three cases that ordinary shape tools conflate: shape-invalid index tensors,
shape-valid but value-invalid index tensors, and value-dependent indices whose
contents are unknown statically.  Stable FX-visible index buffers carry
deterministic integer min/max ranges, so the verifier can report a separate
high-confidence \texttt{index\_bounds} bug when a real index is negative or
exceeds \texttt{input.size(dim)}.  Random FX-folded constants are deliberately not
trusted for value bounds, and PyTorch's real dtype rule is followed exactly:
these operators accept \texttt{int32} or \texttt{int64} indices, while floating or
other integer widths are dtype errors.

\subsection{Padding as a Library-Exact Boundary Rule}

Padding looks simple until the exact PyTorch boundary conditions matter.  The
repository now models \texttt{F.pad} and \texttt{nn.*Pad*d} modules for
\texttt{constant}, \texttt{reflect}, \texttt{replicate}, and \texttt{circular}
modes.  The transfer keeps PyTorch's last-dimension-backward tuple convention,
preserves tuple lengths exactly as supplied, expands integer module padding
according to the concrete module class, and handles negative padding as cropping.
It also enforces the mode-specific runtime constraints: constant mode may pad any
rank as long as the pad tuple fits the tensor rank; reflect/replicate/circular
support only the documented 1D/2D/3D spatial cases; reflect padding must be
smaller than the corresponding concrete input extent; circular padding may not
wrap a dimension more than once; and replication is allowed to pad farther than
the input size when eager PyTorch does.

The tests are differential, not merely algebraic.  They execute real PyTorch for
successful output shapes and for runtime errors, then require source and FX
TensorGuard paths to agree.  This catches exactly the kind of edge case that a
paper-only rule would likely miss: for example, \texttt{nn.ReflectionPad2d((1,2))}
does \emph{not} mean symmetric 2D padding, and a circular pad equal to a
dimension is legal while one larger than the dimension is not.

\subsection{Interpolation and Upsampling}

\texttt{F.interpolate} and \texttt{nn.Upsample} are shape-changing even though
their parameters often look like harmless presentation choices.  \tg{} now
models their literal \texttt{size} and \texttt{scale\_factor} contracts across
1D/2D/3D spatial tensors: scalar sizes broadcast across all spatial axes, tuple
sizes must match spatial rank, and scale factors use PyTorch's integer
\(\lfloor \mathrm{input}\times\mathrm{scale}\rfloor\) output formula, including
fractional downsampling and per-axis tuple scales.  Options such as
\texttt{align\_corners}, \texttt{mode}, \texttt{recompute\_scale\_factor}, and
\texttt{antialias} are captured so literal runtime-invalid combinations are
reported, while dynamic sizes such as \texttt{x.shape[-2:]} conservatively
produce symbolic spatial dimensions instead of false positives.

\subsection{Recurrent Tensor Contracts}

Recurrent modules are another place where a shape-preserving approximation is
wrong.  \tg{} models \texttt{nn.RNN}, \texttt{nn.GRU}, and \texttt{nn.LSTM}
using PyTorch's sequence contract: \texttt{batch\_first} affects only the input
and output layout; hidden states keep \((D\cdot L,N,H)\)-style layout; bidirectionality
doubles the output feature dimension but not each returned hidden direction; and
projected LSTMs use \texttt{proj\_size} for \texttt{output}/\texttt{h\_n} while
\texttt{c\_n} retains \texttt{hidden\_size}.  Multilayer and unbatched cases share
the same contract machinery, so downstream \texttt{Linear}, reshape, and tuple
unpacking checks see the exact returned tensor ranks instead of a last-dimension
heuristic.

Packed sequences are treated as an explicit abstention boundary.  Calls such as
\texttt{pack\_padded\_sequence} produce opaque values; recurrent layers consuming
them do not fabricate dense tensor shapes, and downstream checks remain honest
until a supported dense tensor reappears.  The regression tests include direct
verifier checks for wrong hidden/cell/projection dimensions, a live PyTorch
projected bidirectional LSTM shape comparison, and a packed-sequence false-positive
guard.

\subsection{Examples of Transfer Precision}

Reshape and view rules reason about element products and the \texttt{-1}
inference sentinel.  They reject concrete product mismatches, infer the missing
dimension when possible, and preserve copied symbolic dimensions through
\texttt{x.size(i)} and \texttt{x.shape[i]} aliases.  Broadcasting implements
NumPy/PyTorch right-aligned semantics, including singleton expansion and
symbolic abstention.  Convolution rules compute spatial output sizes from
kernel, stride, padding, dilation, groups, and transpose output padding when
these parameters are static.  Padding rules follow eager PyTorch across
constant, reflect, replicate, and circular modes, including negative cropping and
mode-specific rank/extent errors.  Fold/unfold rules use the same spatial arithmetic
to model batched and unbatched image-column transforms, reject empty sliding
block grids, non-divisible fold channel counts, and output-size/column-count
reconstruction mismatches.  Normalization rules check channel or normalized shape
contracts, while phase annotations make train/eval-sensitive modules diagnostic
instead of invisible.

Attention rules are handled at multiple layers.  SDPA verifies query/key
embedding compatibility and key/value source-length compatibility while
preserving the output value dimension.  Multi-head attention checks packed and
unpacked query/key/value forms, \texttt{kdim}/\texttt{vdim}, batch-first layout,
attention masks, key-padding masks, per-head weight shapes, and valid symbolic
cases.  Tests replay these contracts against real PyTorch modules and functional
calls.

The \texttt{einsum} rule is deliberately parser-backed rather than a rank
heuristic.  It models explicit and implicit outputs, ASCII label ordering,
repeated-label diagonals within an operand, ellipsis blocks before or after named
axes, and output ellipses that appear in non-leading positions.  The public
registry transfer and the verifier use the same parser, and the Z3 encoding
maps output labels after the broadcasted ellipsis positions instead of assuming
that labels always start at dimension zero.

Data-dependent operators are treated honestly.  \texttt{masked\_select},
\texttt{nonzero}, single-argument \texttt{where}, and boolean-mask indexing can
produce extents that depend on runtime tensor values.  \tg{} now first checks
the static part of the contract---mask broadcastability or boolean-mask prefix
shape---and then emits symbolic extents plus explicit \unknown{} reasons instead
of guessing a concrete selected length.  This design is repeated throughout the
codebase: useful exactness where the contract is shape-determined, conservative
abstention where it is not.

The linear-algebra rules illustrate the same standard for multi-output operators.
\texttt{torch.linalg.svd} and \texttt{qr} return named output tuples whose shapes
depend on \(\min(M,N)\), \texttt{full\_matrices}, and QR mode; the \texttt{r}
mode's empty \texttt{Q} output is modeled exactly.  Square-matrix operators
(\texttt{inv}, Cholesky, and eigendecomposition) reject only statically known
rank or trailing-dimension violations.  \texttt{solve} preserves PyTorch's
non-obvious RHS branch: vector-shaped batched RHS tensors do not broadcast like
matrix RHS tensors, and \texttt{left=False} rejects the exact-batch vector branch.
Focused tests compare all of these shapes and failures against the installed
PyTorch dispatcher, while unsupported linalg operators remain heuristic.

The higher-order \texttt{torch.func} rules carry that exactness into
differentiable wrappers.  \texttt{grad} requires a rank-0 scalar body output and
preserves PyTorch's int-versus-sequence \texttt{argnums} return structure.
\texttt{jacrev} and \texttt{jacfwd} map every output pytree leaf with the product
shape ``output plus differentiated input'', \texttt{jvp} validates tangent
shapes, and \texttt{vjp} validates cotangent trees while documenting the pullback
gradient tuple.  Symbolic tangent/cotangent equalities and value-dependent
closures abstain instead of creating false contracts.

Training loss functions are modeled as multi-input contracts rather than unary
shape reducers.  \tg{} distinguishes class-index and probability-target
\texttt{cross\_entropy}, checks NLL class targets, MSE and KL broadcasting,
BCE-with-logits exact target shapes, scalar versus unreduced outputs, class
weights, \texttt{pos\_weight}, and eager-PyTorch dtype failures.  The source and
\texttt{torch.fx} frontends preserve both prediction and target operands for
\texttt{nn.*Loss} modules and functional calls, so malformed training snippets
are rejected before the first optimizer step.

\clearpage

\section{Developer Experience}

The repository presents \tg{} as a tool developers can actually adopt.  The CLI
supports zero-flag verification on common \texttt{nn.Module} files, high-
confidence CI gating, domain flags such as \texttt{--no-device-check}, and
soundness modes.  Diagnostics are source-mapped and explain why an operation is
unsafe, not just that Z3 found a model.  The \texttt{--explain} path reconstructs
the inference chain from inputs through transformations to the failing step.
The \texttt{--fix} path suggests mechanical repairs for supported errors such
as wrong \texttt{Linear.in\_features} or convolution channel counts.

Adoption surfaces include a GitHub Action, SARIF for code scanning, JSON and
JUnit reporters, GitHub PR annotations, pre-commit integration, nox/tox hooks,
a pytest plugin, a watch mode, Jupyter/IPython cell checks, a
\texttt{@tensorguard.checked} decorator, a repository config file, framework
hooks for higher-level trainers, and a public API with type stubs.  The
packaging story includes reproducible wheels, a pinned Z3 range, pipx usage,
Docker, conda-forge metadata, and a documented security policy for untrusted
model files.

The diagnostic design is intentionally evidence-oriented.  A bug report can
carry an expected shape, actual shape, line and column, a counterexample, the
domain that produced the violation, and the chain of operations that made the
state refutable.  When a domain is disabled, the solver encoders for that
domain emit no constraints, so the flag removes solver work instead of merely
filtering the final verdict.  This matters for ablation experiments and for
users who need predictable CI behavior.

\section{Empirical Methodology}

\tg{} ships with a broad evaluation harness.  The artifact is organized around
questions a reviewer or user would ask.

\begin{quote}
Does it catch real bugs?  Does it avoid false positives on clean models?  Does
it agree with PyTorch on operator semantics?  Does it scale?  Which abstract
domains matter?  How does it compare to baselines?  Can the entire claim set be
regenerated from committed scripts?
\end{quote}

The repository answers these questions through several benchlines and
experiments rather than a single benchmark number.  The frozen ground-truth
corpus in \texttt{real\_benchmarks/} contains clean and buggy PyTorch modules
with content hashes and labels.  The GitHub-mined dataset under
\texttt{experiments\_v5/github\_bug\_mining/} records real public issue and PR
signals associated with PyTorch runtime error signatures.  Fuzzing harnesses
generate valid modules and mutation harnesses inject shape, device, and dtype
faults.  Differential conformance tests cross-check transfer functions against
the live torch dispatcher.  Baseline comparison scripts evaluate PyTea and
runtime-tool alternatives.  Dashboards and JSON artifacts track coverage,
precision/recall, latency, ablations, false positives, and statistical
rigor.

\begin{longtable}{|p{0.25\linewidth}|p{0.68\linewidth}|}
\hline
\textbf{Evidence family} & \textbf{Repository artifact and purpose} \\
\hline
Ground-truth corpus &
\texttt{real\_benchmarks/}, \texttt{tests/test\_real\_benchmarks.py}.  Frozen
clean/buggy modules with a manifest, hashes, labels, and loader checks. \\
\hline
GitHub-mined bugs &
\texttt{experiments\_v5/github\_bug\_mining/}.  Offline dataset plus miner for
public shape/device bug candidates from issue and PR history. \\
\hline
Precision/recall &
\texttt{tests/test\_precision\_recall.py},
\texttt{tests/test\_hard\_recall.py}, \texttt{experiments/}.  Tracks caught
bugs, clean-model behavior, and high-confidence safety. \\
\hline
False-positive stress &
\texttt{tests/test\_fp\_stress.py}, \texttt{tests/test\_sound\_mode\_fp.py}.
Exercises clean real-model patterns and sound-mode no-false-positive posture. \\
\hline
Differential dispatch &
\texttt{tests/test\_differential\_dispatcher.py},
\texttt{tests/test\_conformance\_oracle.py},
\texttt{tests/test\_propagator\_contracts\_runtime.py}.  Compares transfer
outputs to live PyTorch behavior. \\
\hline
Fuzzing and mutation &
\texttt{tests/test\_diff\_fuzz.py}, \texttt{tests/test\_neg\_fuzz.py},
\texttt{tests/test\_hypothesis\_module\_ast.py},
\texttt{tests/test\_mutation\_clean\_models.py}.  Searches for disagreements
and injected bugs. \\
\hline
Operator coverage &
\texttt{operator\_confidence\_table.json},
\texttt{tests/test\_operator\_coverage.py},
\texttt{tests/test\_operator\_coverage\_gate.py},
\texttt{tests/test\_real\_model\_operator\_coverage.py}.  Tracks supported
operators, confidence tags, release gates, and frequency-weighted real-model
coverage over torchvision, timm, and HuggingFace traces. \\
\hline
Latency and scaling &
\texttt{tests/test\_cost\_benchmarks.py},
\texttt{tests/test\_latency\_budgets.py},
\texttt{tests/test\_scaling\_study.py}.  Measures model-size and solver-cost
behavior. \\
\hline
Baselines &
\texttt{tests/test\_baselines.py},
\texttt{tests/test\_baseline\_head\_to\_head.py},
\texttt{tests/test\_headtohead\_n15.py}.  Compares against PyTea, runtime
checks, and type-oriented alternatives where applicable. \\
\hline
Ablations &
\texttt{tests/test\_domain\_ablation.py},
\texttt{tests/test\_reduced\_product\_ablation.py},
\texttt{tests/test\_cegar\_depth\_ablation.py},
\texttt{tests/test\_verified\_reduction\_diff.py}.  Measures contribution of
domains and refinement depth. \\
\hline
Dashboards &
\texttt{tests/test\_dashboard.py}, \texttt{tests/test\_build\_dashboard.py},
\texttt{tests/test\_leaderboard.py}.  Regenerates static result views for
reviewers and contributors. \\
\hline
Statistical rigor &
\texttt{tests/test\_statistical\_power.py},
\texttt{tests/test\_significance.py},
\texttt{tests/test\_corrected\_meta\_analysis.py}.  Documents effect sizes,
power, and aggregate evidence. \\
\hline
\end{longtable}

\clearpage

\section{Evaluation Highlights}

The strongest evidence is not one number; it is the agreement between many
small, independently checkable claims.  \tg{}'s benchmark suite exercises real
failure modes: wrong \texttt{Linear} feature counts after flattening, malformed
convolution channels, bad Q/K/V dimensions, attention mask errors, invalid
reshape products, device mismatches between buffers and inputs, dtype
divergences, gradient detach bugs, invalid embedding-bag offsets and
per-sample-weight contracts, invalid sparse compressed layouts and
sparse-dense kernels, named
tensor alignment failures, incompatible loss targets, incompatible distribution
batch/event shapes, export
contracts, quantization observer calibration/qscheme and quant/dequant placement,
mixed-precision/autocast dtype drift, and distributed sharding mistakes.

The conformance tests are especially important.  A static rule is only useful
if it matches the runtime library.  The repository therefore includes tests
that construct real tensors, run the corresponding PyTorch operator, and compare
the resulting shape or runtime exception to \tg{}'s predicted contract.  This
pattern appears for cat, einsum, reshape, broadcasting, convolution, attention,
recurrent modules, fold/unfold, padding, grid sampling, embedding-bag
pooling and jagged-offset contracts, sparse layouts and sparse operators, complex
FFT views, named tensors, distributions, vmap, loss functions, chunk/split
partitioning, and index-value operators.  When PyTorch's
edge-case behavior is surprising, as with \texttt{take\_along\_dim(dim=None)}
flattening before selection or \texttt{topk} rejecting oversized \texttt{k}, the
static rule follows PyTorch rather than a simplified textbook semantics.

The ablation suite makes the multi-domain design measurable.  Shape alone
catches the most common architectural errors, but device and gradient domains
catch bugs that shape-only tools miss.  Dtype reasoning catches mixed-precision
and cast-related failures.  Phase information explains train/eval-sensitive
patterns.  The reduced product is valuable because real failures often cross
domain boundaries: a compile/export path may require a shape, dtype, device,
and layout combination, not merely one axis of correctness.

Latency and scalability scripts ensure the verifier remains useful.  Solver
queries are cached or avoided when a transfer can be decided syntactically.
Constraint simplification performs constant folding, common-subexpression
sharing, divisibility normalization, and short-circuiting.  Per-module
verification can run in parallel, and incremental analysis reuses results when
only part of a model changes.  Timeout and anytime modes return sound partial
results instead of claiming full coverage after a budget expires.

Operator coverage is measured by the operators that actually occur in traced
models, not just by the raw public PyTorch surface.  The latest census traces
torchvision, timm, and no-download HuggingFace models, publishes a before/after
matrix, and pins each newly counted hot operator with real extractor tests:
\texttt{torchvision.ops.stochastic\_depth} is treated as shape-preserving,
functional \texttt{layer\_norm} and \texttt{adaptive\_avg\_pool2d} build
synthetic FX layers, and SDPA maps to the existing attention contract.  Scalar
shape-metadata queries such as \texttt{size}, \texttt{dim}, and integer
\texttt{floordiv} are reported separately instead of being counted as tensor
transfer functions.

\section{Formal and Trust Story}

\tg{} does not ask users to trust only implementation tests.  The repository
contains Lean artifacts for several core rules and proof obligations, including
CEGAR convergence bounds, infeasible-contract abstention, fragment modes,
device/dtype transfer soundness, SMT encoding faithfulness, cross-domain
composition, and operator-specific rules such as cat and reshape families.
The Lean audit checks theorem statements for sorry-free status under the
trusted kernel axioms used by the environment.

The formalization is intentionally scoped.  It does not claim to prove all of
PyTorch.  Instead, it proves the algebraic cores that the implementation relies
on and pairs them with runtime conformance tests for library behavior.  For a
research artifact, this is a practical trust split: Lean validates the abstract
rule, Z3 discharges symbolic conditions in the implementation, and PyTorch
oracle tests validate that the rule corresponds to the library release being
tested.

The soundness contract documents the remaining boundary.  Unsupported or
heuristic operators are not silently certified.  Opaque layers and
data-dependent control flow produce \unknown{} in strict modes.  Symbolic sizes
are preserved instead of guessed.  This is visible in the code.  For chunk/split,
for example, concrete dimensions produce exact per-output sizes and concrete
unpack-count errors; symbolic split dimensions produce fresh symbolic split
axes, preventing false claims about a final chunk size.

\clearpage

\section{Implementation Surfaces}

The system is intentionally broad because tensor bugs are not confined to a
single frontend.  The repository includes modules for source analysis, model
checking, graph compilation, torch integration, export extraction, Dynamo gap
analysis, framework hooks, sparse and complex verification, SDPA and MHA
verification, distribution checks, named tensor checks, vmap checks, quant/export
checks, distributed verification, LoRA verification, training-loop checks,
reporters, SARIF, baselines, dashboards, proof certificates, and API stability.

\begin{longtable}{|p{0.24\linewidth}|p{0.69\linewidth}|}
\hline
\textbf{Surface} & \textbf{Representative capability} \\
\hline
\texttt{src/api.py} &
Public \texttt{analyze}, \texttt{verify\_architecture}, and \texttt{verify\_module}
entry points with soundness modes and domain flags. \\
\hline
\texttt{src/model\_checker.py} &
AST graph extraction, symbolic state propagation, Z3 safety checks, device,
phase, gradient, dtype, recurrent, and shape transitions. \\
\hline
\texttt{src/tensor\_shapes.py} &
Shape algebra, transfer helpers, AST shape analyzer, reshape/broadcast/cat and
chunk/split shape computation. \\
\hline
\texttt{src/fx\_extractor.py} and \texttt{src/export\_extractor.py} &
Tracing and export-oriented frontends for real modules and deployment flows. \\
\hline
\texttt{src/torch\_integration.py} &
Guarded compile/export/package pre-passes that surface \tg{} violations before
compiler or exporter failures. \\
\hline
\texttt{src/mha\_verify.py}, \texttt{src/sdpa\_verify.py} &
Attention-family contracts for q/k/v, masks, head counts, and output shapes. \\
\hline
\texttt{src/embedding\_bag\_verify.py} &
PyTorch and TorchRec embedding-bag contracts for offsets, per-sample weights,
pooling modes, index bounds, and jagged/ragged abstention. \\
\hline
\texttt{src/sparse\_verify.py}, \texttt{src/complex\_verify.py} &
Sparse layout, sparse-dense operator, conversion, and complex/FFT shape-dtype
contracts. \\
\hline
\texttt{src/distributions\_verify.py}, \texttt{src/named\_tensor\_verify.py},
\texttt{src/vmap\_verify.py}, \texttt{src/func\_autodiff\_verify.py},
\texttt{src/loss\_verify.py}, \texttt{src/torchvision\_v2\_verify.py} &
High-value library contracts beyond ordinary dense layers, including image
preprocessing transforms whose tensor shapes feed the model. \\
\hline
\texttt{src/quant\_export\_checks.py}, \texttt{src/mixed\_precision\_verify.py},
\texttt{src/distributed\_verification.py}, \texttt{src/lora\_verification.py} &
Deployment, mixed-precision, distributed training, and adapter-compatibility
gates. \\
\hline
\texttt{src/reporters.py}, \texttt{src/sarif\_codescan.py},
\texttt{src/github\_action.py} &
Machine-readable and platform-native reporting. \\
\hline
\texttt{lean/TensorGuard/*.lean} &
Machine-checked statements for selected transfer rules, mode boundaries,
encoding faithfulness, and CEGAR reasoning. \\
\hline
\end{longtable}

\section{Artifact Reproducibility}

The artifact is designed to be regenerated.  The \texttt{make reproduce}
harness and reproducibility scripts update headline results, numeric claim
audits, benchmark deltas, scaling curves, CEGAR measurements, and paper evidence
indices.  Tests pin committed JSON artifacts against freshly generated outputs.
The benchmark corpus has content hashes, datasheets, citation metadata, and
loader checks.  The operator confidence table is generated from code.  The
dashboard is regenerated from committed results.  The public docs explain what
the verifier can and cannot prove.

Reproducibility is paired with hygiene.  Third-party source trees are
quarantined or removed unless license-compatible.  The security policy and
threat model acknowledge that untrusted model files are parsed and should not
be executed dynamically.  Safe loading paths use AST/fx-style extraction rather
than arbitrary imports where possible.  Baselines and suppressions let teams
adopt the tool incrementally without hiding new regressions.

\clearpage

\section{Limitations}

\tg{} is intentionally conservative.  It is not a complete theorem prover for
all Python or all PyTorch.  It handles a documented verifiable fragment and
falls back to \unknown{} outside that fragment in strict modes.  Data-dependent
shape operators, arbitrary Python side effects, dynamic module mutation,
runtime-generated code, and library internals that cannot be traced or modeled
may require abstention.  Symbolic arithmetic is limited to decidable fragments
and practical solver budgets.  Some integrations verify preconditions for
export, compile, serving, quantization, or distributed execution rather than
replacing those systems' own validators.

These limitations are design choices, not hidden caveats.  The repository keeps
an honest limitations document, a formal fragment specification, confidence
tags for operators, and tests that assert unsupported constructs become
\unknown{} rather than \safe{}.  The most important invariant is not that every
program is proved; it is that a strict safe verdict means what it says.

\section{Why the Artifact Matters}

For a tool to be credible at PLDI, OOPSLA, ISSTA, or a systems-for-ML venue, it
must connect a formal idea to the messiness of real code.  \tg{} does that in
three ways.  First, it makes refinement-style tensor verification usable from
normal PyTorch files with zero annotations.  Second, it exposes enough product
surface--CLI, CI, SARIF, Jupyter, decorators, framework hooks, export gates,
Docker, conda, pipx, docs--that adoption can be evaluated rather than imagined.
Third, it ships an artifact that reviewers can run, audit, and extend.

The codebase also provides a platform for future work.  New operator rules can
be added with confidence tags, conformance tests, and proof footprints.  New
benchmarks can enter the frozen corpus through provenance and hashing.  New
frontends can lower into the same graph representation.  New proof-carrying
certificates can reuse the existing safety certificate and replay infrastructure.
The next natural milestones are larger blind real-bug corpora, broader
deployment galleries, independently checkable certificates, and a community
leaderboard that makes tensor-verifier progress measurable.

\section{Conclusion}

\tg{} demonstrates that static verification for PyTorch tensor programs can be
both principled and practical.  The repository combines a refinement-type view
of tensors, a multi-domain abstract state, Z3-backed checks, CEGAR refinement,
soundness modes, precise operator transfer functions, real frontends, formal
artifacts, and an unusually broad reproducibility harness.  The latest
gather/scatter and padding work illustrates the project's standard: match
PyTorch's exact edge-case semantics, propagate precision to downstream model
checks, abstain on symbolic or value-dependent uncertainty, and pin the behavior
with live-library tests.

The result is a verifier that catches high-value architecture, device, dtype,
phase, gradient, export, and deployment bugs before expensive training or
opaque compiler failures.  It is a research artifact with a clear trust
boundary and a developer tool with practical adoption paths--the combination
needed for a best-paper-quality static analysis system in modern ML software.

\clearpage

\appendix
\section{Capability Matrix}

\begin{longtable}{|p{0.27\linewidth}|p{0.64\linewidth}|}
\hline
\textbf{Capability} & \textbf{What the repository demonstrates} \\
\hline
Zero-annotation source verification &
Inference from constructors, layer definitions, tensor factories, asserts,
shape aliases, and dataflow without requiring user type annotations. \\
\hline
Shape algebra &
Matmul, reshape/view, flatten, broadcasting, cat, stack-family aliases, chunk, split,
padding, transpose, permute, squeeze, unsqueeze, reductions, indexing,
convolution, pooling, fold/unfold, embedding, and attention rules. \\
\hline
Device reasoning &
Propagation through tensor creation, module calls, \texttt{to}/\texttt{cuda}/
\texttt{cpu}, buffers, and cross-device operations. \\
\hline
Dtype reasoning &
Promotion/cast checks, explicit autocast backend/dtype contexts, AMP allowlists,
GradScaler requirements, parameter/input dtype matching, quantization/export
dtype gates, and dtype-related runtime-error prevention. \\
\hline
Phase reasoning &
Train/eval tagging and phase-sensitive diagnostics for modules such as
BatchNorm and Dropout. \\
\hline
Gradient reasoning &
Detection of detached or broken gradient flow and training-loop checks around
optimizers, parameters, and checkpoint-like patterns. \\
\hline
Training losses &
Target-shape, reduction, dtype, class-weight, and \texttt{pos\_weight}
contracts for cross entropy, NLL, MSE, BCE-with-logits, and KL losses. \\
\hline
Stride and permutation &
Layout and viewability support that keeps reshape/permute reasoning from
collapsing to rank-only approximations. \\
\hline
Attention &
SDPA and MHA contracts, q/k/v dimensions, kdim/vdim, masks, key padding masks,
batch-first layout, and grouped head variants. \\
\hline
Embedding bags &
\texttt{nn.EmbeddingBag}, \texttt{F.embedding\_bag}, table dimensions, offsets,
\texttt{include\_last\_offset}, per-sample weights, pooling modes, index
bounds, and TorchRec-style jagged abstention. \\
\hline
Sparse tensors &
COO/CSR/CSC/BSR/BSC shape and blocksize invariants, sparse-dense
\texttt{mm}/\texttt{addmm}, CSR \texttt{sampled\_addmm}, sparse softmax,
coalesce, \texttt{to\_dense}, and layout conversions. \\
\hline
Torchvision v2 transforms &
Source-level tensor contracts for resize, crops, padding, normalization,
flips/color/dtype transforms, and PIL-only abstention. \\
\hline
Named tensors &
\texttt{refine\_names} and \texttt{align\_to} alignment checks. \\
\hline
Complex and FFT &
\texttt{view\_as\_real}, \texttt{view\_as\_complex}, FFT shape/dtype contracts. \\
\hline
Distributions &
Batch shape, event shape, sample shape, and log-probability compatibility for
probabilistic code. \\
\hline
vmap/functorch &
Batched dimension transfer and out-dim contract checks. \\
\hline
Export and compile &
Pre-pass verification for \texttt{torch.compile}, \texttt{torch.export}, ONNX,
AOTInductor packaging, and deployment gates, including package-time layout,
dtype, device, dynamic-guard, and lowering checks. \\
\hline
Distributed and adapters &
DDP/FSDP/DTensor/tensor-parallel checks, optimizer-state checks, checkpoint
schema checks, LoRA/PEFT adapter compatibility. \\
\hline
Reports and integrations &
CLI, JSON, JUnit, SARIF, GitHub annotations, GitHub Action, pre-commit, pytest
plugin, nox/tox, watch mode, Jupyter, decorators, and framework hooks. \\
\hline
Formal artifacts &
Lean proof files, axiom audits, SMT encoding faithfulness checks, and proof
certificate infrastructure. \\
\hline
Reproducibility &
Frozen corpora, generated dashboards, numeric claim audits, benchmark scripts,
hash manifests, datasheets, artifact badges, Docker/conda/source modes. \\
\hline
\end{longtable}

\section{Benchline and Experiment Inventory}

\begin{longtable}{|p{0.30\linewidth}|p{0.61\linewidth}|}
\hline
\textbf{Artifact} & \textbf{Role in the evaluation story} \\
\hline
\texttt{tests/test\_stack\_family.py} &
Live PyTorch conformance for \texttt{torch.stack}, \texttt{hstack},
\texttt{vstack}, \texttt{dstack}, \texttt{column\_stack}, and
\texttt{row\_stack}: singleton TensorLists, negative dims, scalar/1D/2D/3D
rank promotions, source/FX extraction, and static error parity. \\
\hline
\texttt{tests/test\_split\_chunk\_precise.py} &
Live PyTorch cross-checks and verifier regressions for exact chunk/split
partition rules, empty chunks, negative dims, function/method forms, and
downstream Linear/cat behavior. \\
\hline
\texttt{tests/test\_mha\_verify.py} &
Multi-head attention conformance against real \texttt{nn.MultiheadAttention}
and functional attention calls. \\
\hline
\texttt{tests/test\_einsum\_precise.py} &
Equation parser and output-shape tests for einsum, including repeated-label
diagonals, ellipsis placement, ASCII output ordering, registry-transfer parity,
and Z3 output-position constraints. \\
\hline
\texttt{tests/test\_index\_value\_ops\_precise.py},
\texttt{tests/test\_gather\_scatter\_index\_bounds.py} &
Live PyTorch conformance for \texttt{take\_along\_dim}, \texttt{argsort},
\texttt{sort}, \texttt{topk}, \texttt{kthvalue}, \texttt{argmax}, and
\texttt{argmin}, plus gather/scatter/index-select index dtype and static
value-bound errors. \\
\hline
\texttt{tests/test\_reshape\_precise.py},
\texttt{tests/test\_reshape\_divisibility\_opt.py} &
Product reasoning, \texttt{-1} inference, and optimized divisibility
constraints for reshapes/views. \\
\hline
\texttt{tests/test\_broadcast\_expand\_precise.py},
\texttt{tests/test\_broadcast\_correspondence.py} &
Broadcast and expand semantics, including singleton expansion and torch
correspondence. \\
\hline
\texttt{tests/test\_conv1d\_rule.py},
\texttt{tests/test\_conv2d\_rule.py},
\texttt{tests/test\_conv3d\_rule.py},
\texttt{tests/test\_pool2d\_rule.py},
\texttt{tests/test\_convtranspose\_rule.py} &
Convolution and pooling transfer functions over spatial arithmetic. \\
\hline
\texttt{tests/test\_unfold\_fold\_spatial.py} &
Live PyTorch conformance for \texttt{nn.Unfold}, \texttt{nn.Fold},
\texttt{F.unfold}, and \texttt{F.fold}, including batched/unbatched ranks,
dilation/padding/stride block counts, empty grids, non-divisible fold channels,
and impossible fold column-count/output-size reconstructions. \\
\hline
\texttt{tests/test\_pad\_semantics.py} &
Live PyTorch conformance for \texttt{F.pad} and \texttt{nn.*Pad*d}: exact
constant/reflect/replicate/circular output shapes, source and FX parity,
negative cropping, non-reinterpreted tuple lengths, 3D circular modules, and
runtime-error parity for invalid ranks, reflect extents, and circular over-wraps. \\
\hline
\texttt{tests/test\_interpolate\_precise.py} &
Live PyTorch conformance for \texttt{F.interpolate} and \texttt{nn.Upsample}:
literal \texttt{size} vs. \texttt{scale\_factor}, fractional and tuple scales,
rank-3/4/5 formulas, source and FX parity, dynamic-size abstention, downstream
shape flow, and runtime-error parity for invalid modes/options. \\
\hline
\texttt{tests/test\_dtype\_precise.py},
\texttt{tests/test\_dtype\_promote\_chain.py},
\texttt{tests/test\_bool\_dtype\_soundness.py},
\texttt{tests/test\_mixed\_precision\_verify.py} &
Dtype propagation, promotion, boolean/integer safety, and live PyTorch
CPU-autocast parity for mixed-precision policy gates. \\
\hline
\texttt{tests/test\_loss\_verify.py} &
Loss-function conformance for target shape, reduction, dtype, class weights,
\texttt{pos\_weight}, source snippets, functional calls, FX modules, and eager
PyTorch runtime-error parity. \\
\hline
\texttt{tests/test\_torchvision\_v2\_verify.py} &
Live torchvision v2 conformance for resize/crop/pad/normalize/flips, source
\texttt{Compose} and inline constructor extraction, runtime-error parity, and
PIL-only symbolic abstention. \\
\hline
\texttt{tests/test\_device\_propagation\_precise.py},
\texttt{tests/test\_device\_placement\_transfer.py} &
Device movement and cross-device error detection. \\
\hline
\texttt{tests/test\_grad\_flow\_transfer.py},
\texttt{tests/test\_training\_loop\_checks.py} &
Gradient and training-loop contracts beyond forward shape checking. \\
\hline
\texttt{tests/test\_onnx\_export\_gate.py},
\texttt{tests/test\_export\_aot\_gate.py},
\texttt{tests/test\_quant\_export\_checks.py},
\texttt{tests/test\_quantization\_verify.py} &
Deployment-facing gates for export, \texttt{torch.export.Dim} ranges and
relations, opsets, AOT package input/layout/dtype/device/lowering contracts, and
source/eager/FX quantization. \\
\hline
\texttt{tests/test\_distributed\_verification.py},
\texttt{tests/test\_tensor\_parallel\_checks.py} &
Distributed and tensor-parallel shape/device contracts. \\
\hline
\texttt{tests/test\_lora\_verification.py} &
Adapter rank and target-module compatibility checks. \\
\hline
\texttt{tests/test\_security.py},
\texttt{tests/test\_threats\_to\_validity.py} &
Security policy and evaluation-threat documentation. \\
\hline
\texttt{tests/test\_docs\_getting\_started.py},
\texttt{tests/test\_docs\_limitations.py} &
Executable documentation checks for onboarding and limitations. \\
\hline
\texttt{tests/test\_paper\_evidence\_index.py},
\texttt{tests/test\_paper\_capsule\_wiring.py} &
Paper/artifact wiring so claims can be tied to scripts and committed evidence. \\
\hline
\end{longtable}

\section{Transfer Rule Ledger}

This appendix summarizes representative transfer rules as contracts rather than
as an operator checklist.  The purpose is to make clear which parts of the
system are exact, which are sound but conservative, and which intentionally
abstain.

\begin{longtable}{|p{0.22\linewidth}|p{0.31\linewidth}|p{0.36\linewidth}|}
\hline
\textbf{Family} & \textbf{Static contract} & \textbf{Evidence path} \\
\hline
Tensor factories &
\texttt{zeros}, \texttt{ones}, \texttt{rand}, \texttt{randn}, \texttt{empty},
\texttt{full}, and integer factories produce shapes determined by static size
arguments, independent of random values. &
Factory transfer tests and RNG shape tests; output shapes are propagated even
when tensor contents are unknown. \\
\hline
Linear layers &
The last input dimension must equal \texttt{in\_features}; output replaces the
last dimension with \texttt{out\_features}. &
Unit tests for clean and buggy linear models, autofix tests, and end-to-end
module verification. \\
\hline
Recurrent modules &
\texttt{RNN}, \texttt{GRU}, and \texttt{LSTM} check input feature size, preserve
sequence/batch layout, account for layers and directions, distinguish LSTM
\texttt{h\_n} from \texttt{c\_n}, support projections, and abstain on packed
sequences. &
Focused recurrent tests plus a live PyTorch projected bidirectional LSTM
shape oracle and packed-sequence abstention regression. \\
\hline
Matrix multiplication &
Contracting dimensions must match; batched prefixes are propagated or
broadcast where supported. &
Shape arithmetic tests, Z3 symbolic checks, model-checker violations, and
baseline bug cases. \\
\hline
Convolution &
Input rank, channels, groups, kernel, stride, padding, dilation, and output
padding determine channel and spatial output shapes. &
Conv1d/2d/3d, transpose convolution, pooling, and pixel-shuffle regression
tests. \\
\hline
Fold/unfold &
\texttt{unfold} maps image tensors to image-column tensors with exact
kernel/dilation/padding/stride block counts; \texttt{fold} inverts the shape only
when channel divisibility and the output-size-implied sliding-block product
match the input columns. &
Live PyTorch differential tests for module and functional forms, including
runtime-error parity for empty grids and impossible reconstruction counts. \\
\hline
Padding &
\texttt{F.pad} and \texttt{nn.*Pad*d} apply tuple pairs from the last dimension
backward; constant mode accepts any compatible rank, while reflect, replicate,
and circular modes enforce PyTorch's spatial-rank cases and concrete extent
constraints. Negative padding crops where PyTorch permits it. &
Live PyTorch differential tests for functional and module forms across source
and FX, including invalid-rank, reflect-bound, circular-wrap, and large
replicate-padding cases. \\
\hline
Interpolation &
\texttt{F.interpolate} and \texttt{nn.Upsample} preserve batch/channel axes and
compute spatial axes from literal \texttt{size} or
\texttt{scale\_factor}; dynamic size specs abstain soundly. &
Live PyTorch differential tests for fractional scales, tuple scales, 1D/2D/3D
spatial ranks, source/FX parity, invalid option combinations, and downstream
linear-shape checking. \\
\hline
Reshape/view &
Element product is preserved; \texttt{-1} is inferred when uniquely determined;
copied dimensions from \texttt{x.size(i)} and \texttt{x.shape[i]} retain
aliases. &
Precise reshape tests, divisibility optimization tests, Lean-adjacent proof
footprints, and real model regressions. \\
\hline
Flatten/unflatten &
Flatten multiplies a contiguous dimension range when concrete and otherwise
preserves rank with a symbolic product; unflatten checks product compatibility. &
Flatten, unflatten, and reshape-chain tests. \\
\hline
Broadcast and expand &
Right-aligned dimensions must be equal or singleton; expand may only grow
singleton dimensions, and \texttt{-1} preserves an existing dimension where
PyTorch allows it. &
Broadcast correspondence, expand precision, and elementwise-operation tests. \\
\hline
Cat &
All non-concatenation dimensions must match; concatenation dimension sums
concrete sizes or becomes symbolic. Negative dims are normalized. &
Lean cat rule, live PyTorch cat tests, analyzer and model-checker cat
regressions. \\
\hline
Stack &
Base \texttt{stack} requires equal ranks and static-equal concrete dimensions,
then inserts a new axis; \texttt{hstack}/\texttt{vstack}/\texttt{dstack}/
\texttt{column\_stack}/\texttt{row\_stack} apply PyTorch's exact scalar,
1D, 2D, and 3D rank promotions before concatenating along the library-defined
axis. &
Live PyTorch differential tests in \texttt{tests/test\_stack\_family.py},
including source and FX extraction plus runtime-error parity. \\
\hline
Chunk/split &
Exact PyTorch partition sizes for concrete axes; concrete unpack-count errors;
fresh-symbolic split axis when output count or size depends on symbols. &
Live torch 2.9.1 cross-checks in \texttt{tests/test\_split\_chunk\_precise.py}
and downstream Linear/cat model checks. \\
\hline
Squeeze/unsqueeze &
Squeeze removes singleton axes when statically known; unsqueeze inserts a
singleton axis with negative-dim normalization. &
Tensor-shape analyzer tests and model-checker step propagation. \\
\hline
Transpose/permute &
Axis swaps and permutations reorder shape dimensions; identity permutes can be
diagnosed as suspicious no-ops. &
Transpose tests, permutation theory tests, identity-permute diagnostics. \\
\hline
Reductions &
Dim reductions remove an axis or replace it with 1 under \texttt{keepdim};
global reductions produce scalar shapes. &
Reduction tests for mean, sum, min, max, variance, standard deviation, and
arg reductions. \\
\hline
Indexing and gather &
Gather and \texttt{take\_along\_dim} return index-derived shapes
(\texttt{dim=None} flattens); index-select replaces the selected dimension;
scatter preserves the input shape; narrow replaces a dimension with a known
length; sort/topk/kthvalue preserve value dtypes while producing \texttt{int64}
indices; stable gather/scatter/index-select index buffers check dtype and
static bounds while value-dependent indices abstain. &
Indexing/gather tests, index-value precise tests, gather/scatter bounds tests,
and extended operator tests. \\
\hline
Masked and value-dependent ops &
Masked select, nonzero, one-argument where, and boolean masks check static mask
contracts, then use symbolic selected lengths and explicit unknown reasons for
value-dependent extents. &
Boolean-mask tests against live PyTorch plus soundness-boundary tests for
data-dependent ranks. \\
\hline
Einsum &
Equation parser handles explicit and implicit outputs, ellipses, broadcasting,
ASCII output order, and repeated-label diagonal constraints where supported. &
Einsum theory tests, precise einsum tests, and randomized torch differential
batteries. \\
\hline
Attention &
Q/K/V ranks, embedding dimensions, source lengths, batch layout, masks, and
head-related contracts are checked for SDPA and MHA. &
MHA verification tests against real PyTorch modules and functional calls;
SDPA precise tests. \\
\hline
Normalization &
BatchNorm, LayerNorm, GroupNorm, and InstanceNorm preserve shapes while checking
channel or normalized-shape constraints when static. &
Normalization rule tests and phase-flow tests. \\
\hline
Sparse layouts &
Compressed and coordinate layouts carry index/value shape invariants and
blocksize constraints. &
Sparse verification tests across COO, CSR, CSC, BSR, and BSC. \\
\hline
Complex and FFT &
Complex view conversions enforce final-dimension and dtype contracts; FFT
families compute output dimensions such as real FFT half-spectrum sizes. &
Complex verification tests, FFT tests, and operator confidence tags. \\
\hline
Linear algebra &
\texttt{solve}, inverse, Cholesky, SVD, eigendecomposition, and QR enforce
rank, square-matrix, batch-broadcast, mode, and named output-tuple contracts. &
Linalg differential tests against real PyTorch plus sound per-op confidence
tags. \\
\hline
Named tensors &
Name refinement and alignment require axis-name consistency and legal reorder
operations. &
Named tensor verification tests. \\
\hline
Distribution shapes &
Batch, event, sample, and log-probability shapes must be mutually compatible. &
Distribution verification tests and probabilistic-model examples. \\
\hline
vmap/functorch and \texttt{torch.func} autodiff &
Batched dimensions are inserted, removed, or moved according to \texttt{in\_dims}
and \texttt{out\_dims}; autodiff wrappers produce gradient/Jacobian/JVP/VJP
pytrees with PyTorch-faithful \texttt{argnums}, tangent, and cotangent shape
contracts; impossible concrete contracts are rejected. &
vmap verification tests plus \texttt{tests/test\_func\_autodiff\_verify.py}
live dispatcher checks. \\
\hline
Export gates &
Export, ONNX, AOTInductor, quantization, and mixed-precision paths check static
preconditions before invoking downstream compilers; dynamic
\texttt{torch.export.Dim} ranges are reconciled with TensorGuard input-shape
symbols before tracing, and AOT packages additionally reject invalid example
layouts, devices, dtypes, missing dynamic guards, and denied ATen lowerings
before the packager writes an artifact. &
ONNX, AOT, quantization, mixed-precision, and guarded compile/export tests. \\
\hline
Distributed and adapter gates &
Sharding, tensor parallelism, optimizer state, checkpoints, and LoRA adapters
have shape/dtype/rank contracts layered over base verification. &
Distributed verification, tensor-parallel, training-loop, checkpoint, and LoRA
tests. \\
\hline
\end{longtable}

\section{Experiment Protocol Notes}

The empirical suite is intentionally redundant.  Each protocol checks a
different failure mode in the claim stack, and overlap is useful: a bug caught
by a real-bug corpus, a mutation benchmark, and a live torch conformance test is
less likely to be an artifact of one benchmark's construction.

\subsection{Real-bug protocol}

Real bugs are valuable only when provenance is preserved.  The repository's
corpus entries carry labels, hashes, and loader checks so that a benchmark case
cannot silently drift.  GitHub-mined cases are derived from public issue and PR
signals, especially runtime error messages that indicate shape or device
failures.  The intent is not to treat every mined candidate as a theorem; it is
to create a large, auditable pool that can be stratified, manually labeled,
filtered into blind splits, and reused by future tools.

\subsection{Clean-model protocol}

False positives are more damaging than missed bugs for a CI gate.  Clean-model
tests therefore exercise patterns that static tools often mishandle: symbolic
batch sizes, shape aliases, helper functions, container modules, attention
wrappers, non-divisible but valid partitions, singleton broadcasting, empty
sections, and deployment preconditions that should abstain rather than refute.
The strict soundness mode is evaluated on these clean cases to maintain the
promise that a trusted safe verdict is not an accident of permissive defaults.

\subsection{Differential protocol}

For operator semantics, PyTorch is the executable specification.  Differential
tests generate tensors, run PyTorch, and compare the observed shape or exception
to \tg{}'s transfer function.  This protocol is especially useful for edge cases
whose behavior is not obvious from memory: \texttt{torch.chunk} returning fewer
than requested chunks, zero-size split sections, broadcasting of attention
masks, padding tuple/rank corner cases, compressed sparse block constraints, and
named-axis alignment.  The same oracle style now covers embedding-bag corner
cases such as 1D inputs requiring offsets, 2D inputs forbidding offsets,
\texttt{include\_last\_offset} changing bag count, and per-sample weights being
legal only in \texttt{mode="sum"}.  It also covers torchvision v2 image
transforms, where \texttt{CenterCrop} accepts rank-2 tensors, \texttt{Resize}
and \texttt{Pad} require tensor-image rank on the tested runtime, and
\texttt{Normalize} can broadcast a singleton channel dimension.
Recent linalg tests add another such case: QR's \texttt{mode="r"} empty
\texttt{Q}, SVD full-vs-reduced output tuples, and \texttt{solve}'s vector/matrix
RHS disambiguation are all dispatcher-visible but easy to misremember.

\subsection{Mutation protocol}

Mutation tests start from clean models and introduce controlled faults:
wrong feature counts, wrong channels, bad reshape products, device moves,
dtype casts, gradient detaches, invalid embedding-bag offsets/weights, invalid
sparse metadata, and export-incompatible
layouts.  A mutation is useful only if it represents a runtime or semantic
failure that real code could contain.  The harness therefore records the mutant,
the expected outcome, and the evidence that the mutation is admitted by the
source syntax but rejected by the verifier or runtime.

\subsection{Ablation protocol}

Ablations isolate the contribution of each abstract domain and each refinement
component.  Shape-only, device-only, dtype-only, gradient-only, independent
domains, reduced products, CEGAR depths, operator coverage levels, solver
timeouts, and frontend variants can be compared.  This is the right protocol
for answering whether the tool's complexity is justified.  If a domain never
catches a unique bug, it should either be simplified or documented as diagnostic.

\subsection{Latency protocol}

Static verification must run before it can help.  The latency protocol measures
small, medium, and large models; import-time, install-time, and analysis-time
costs; solver calls avoided by syntactic transfer functions; incremental
reverification under diffs; parallel per-module verification; and timeout
behavior.  The important metric is not only mean runtime but the tail behavior
under symbolic constraints and unsupported regions.

\begin{longtable}{|p{0.25\linewidth}|p{0.31\linewidth}|p{0.33\linewidth}|}
\hline
\textbf{Protocol} & \textbf{Primary risk controlled} & \textbf{Representative
assertion} \\
\hline
Real-bug corpus & Overfitting to toy bugs & Frozen labels and hashes reproduce
expected verdicts. \\
\hline
Clean-model stress & CI-hostile false positives & Strict mode does not refute
valid real-model patterns. \\
\hline
Differential torch oracle & Incorrect operator semantics & Static transfer
matches PyTorch shape or exception on generated cases. \\
\hline
Mutation benchmark & Weak recall on realistic edits & Injected faults are
detected at the intended operation site. \\
\hline
Fuzzing & Hidden disagreement surfaces & Minimized disagreements become
regression tests. \\
\hline
Baseline comparison & Unclear value over existing tools & Precision, recall,
latency, and time-to-detect are compared against external alternatives. \\
\hline
Domain ablation & Unjustified abstract domains & Removing a domain changes
measurable recall, precision, diagnostics, or runtime. \\
\hline
CEGAR depth ablation & Refinement over- or under-use & More refinement improves
specific cases until convergence without unbounded runtime. \\
\hline
Operator coverage gate & Silent regression in transfer surface & Coverage and
confidence tables remain in sync with code. \\
\hline
Dashboard regeneration & Stale paper claims & Committed result artifacts match
freshly regenerated outputs. \\
\hline
Statistical checks & Underpowered headline claims & Effect-size and sample-size
artifacts document uncertainty. \\
\hline
Artifact packaging & Reproducibility friction & Docker, conda, source, and
pipx paths can run the advertised checks. \\
\hline
\end{longtable}

\section{Adoption Scenarios}

\subsection{Local model authoring}

A researcher writing a new module can run \texttt{tensorguard verify model.py}
before starting training.  The verifier infers shapes from common constructors,
input hints, docstrings, and structural layer usage.  If it finds a bad
\texttt{Linear} feature count or attention mask, it reports a source-mapped
diagnostic.  If it reaches unsupported dynamic behavior, strict mode reports
\unknown{} with an explanation.  This keeps the feedback loop close to the
edit, not to the first expensive batch.

\subsection{Continuous integration}

In CI, teams can run high-confidence or sound mode on changed model files.  The
GitHub Action can annotate pull requests and SARIF can surface findings in code
scanning.  Baseline/suppression support lets a legacy repository adopt the tool
incrementally: existing issues can be acknowledged while new regressions fail
the build.  Domain flags make ablation and staged rollout practical; for
example, a team may initially enforce shape and dtype while logging phase
diagnostics.

\subsection{Notebook exploration}

Jupyter/IPython integration checks a model cell when it is defined.  This mode
is intentionally lightweight.  A notebook user benefits from early shape and
device feedback without building a package or writing a config file.  Because
notebooks often contain dynamic or partially executed state, honest abstention
is important: the tool should explain what it could not prove rather than
pretend that notebook context is a static module.

\subsection{Framework integration}

Framework authors can call the public API before launching training, compiling
a model, exporting an artifact, or serving a checkpoint.  Hooks for Lightning,
HuggingFace-style trainers, accelerate-like launchers, MLflow, W\&B, and CI
systems allow \tg{} to act as a precondition gate.  The value is not merely a
shape checker; it is a common safety layer that can speak JSON, SARIF, JUnit,
or native annotations.

\subsection{Deployment gate}

Deployment introduces contracts that ordinary training does not: ONNX opset
availability, \texttt{torch.export} dynamic-shape \texttt{Dim} ranges and
relations, AOTInductor input layouts, dtype/device policies, preserved dynamic
guards and package-lowering denylists, quantization observer calibration,
activation qscheme placement, eager/FX quant-dequant boundaries,
mixed-precision dtype boundaries, AMP backend allowlists, fp16 GradScaler
requirements, TorchServe or FastAPI request schemas, GGUF or llama.cpp
checkpoint shapes, CUDA graph capture assumptions, and distributed shard
boundaries.  \tg{}'s deployment modules check these contracts before the
downstream system emits a harder-to-debug compiler, exporter, or serving error.

\subsection{Research artifact review}

A reviewer can start from the soundness contract, run the new tests for any
claimed operator family, inspect the evidence index, regenerate dashboards, and
audit Lean files.  This is a different posture from a one-off prototype: the
paper claims are connected to code paths, tests, and artifacts, and unsupported
regions are part of the documented result.

\clearpage

\section{Design Decisions Worth Preserving}

\textbf{Abstention is a feature.}  A static tensor verifier that guesses past
unknown code can appear impressive on demos while being unusable in CI.  \tg{}
instead treats \unknown{} as a first-class outcome.  This enables strict
adoption: a user can decide whether an unknown is acceptable for their workflow.

\textbf{Concrete bugs should be concrete.}  When a shape is known, the diagnostic
should name the operation and the mismatched dimension.  ``Linear expects last
dim=4, got 2'' is more useful than ``solver found SAT.''  Counterexamples are
still valuable, but they should support the explanation rather than replace it.

\textbf{Operator rules must match the live library.}  PyTorch changes, and some
operator behavior is subtle.  Transfer functions should be cross-checked against
the installed PyTorch version wherever possible.  Padding and chunk/split are
model examples: intuitive symmetric padding or equal partitioning would be wrong,
so the rules follow observed torch behavior.

\textbf{Integration surfaces are part of correctness.}  A verifier that catches
bugs only through an internal Python function is not enough.  CI, SARIF,
pre-commit, notebooks, watch mode, export gates, and framework hooks determine
whether the analysis affects real development.

\textbf{Evidence should regenerate.}  The repository's JSON artifacts, dashboards,
coverage tables, and paper evidence files exist so claims can be checked.  A
best-paper-grade artifact must make it easy to see when a number or guarantee is
stale.

\textbf{Formalization should target the trust core.}  Proving all of PyTorch is
not the goal.  Proving the algebra that the implementation relies on, auditing
axioms, and pairing proofs with conformance tests creates a tractable and useful
trust story.

\section{Diagnostic Taxonomy}

The user-facing product of the verifier is a diagnostic, not a theorem object.
The following taxonomy keeps diagnostics comparable across domains and output
formats.

\begin{longtable}{|p{0.24\linewidth}|p{0.31\linewidth}|p{0.34\linewidth}|}
\hline
\textbf{Diagnostic class} & \textbf{Typical trigger} & \textbf{Useful response} \\
\hline
Dimension mismatch & A contract expects two static dimensions to be equal and
they are not. & Report both dimensions, the axis, and the operation that imposed
the equality. \\
\hline
Rank mismatch & An operator requires a tensor rank or rank range that the input
does not satisfy. & Report expected rank, observed rank, and whether a reshape,
unsqueeze, or batch dimension is likely involved. \\
\hline
Product mismatch & A reshape, view, flatten, or unflatten changes the number of
elements. & Report the source product, target product, and the unresolved
symbolic factors if any. \\
\hline
Partition mismatch & Split, chunk, unbind, or unpacking returns a different
number of values than the program expects. & For concrete cases, report the
exact returned count; for symbolic cases, abstain with a dependency explanation. \\
\hline
Broadcast failure & Elementwise or mask shapes cannot be right-aligned under
singleton expansion. & Report the first incompatible aligned axis and both input
shapes. \\
\hline
Device conflict & Inputs to an operation reside on incompatible devices or an
implicit host/device transfer would violate the selected policy. & Report the
producer operation for each device and any missing explicit move. \\
\hline
Dtype conflict & Mixed dtypes are unsupported, lossy, or inconsistent with an
operator's contract. & Report promoted or required dtype and the tensor whose
dtype blocks the operation. \\
\hline
Gradient break & A training path detaches, casts, mutates, or guards tensors in a
way that prevents intended gradient flow. & Report the first operation that
breaks differentiability and the downstream loss or parameter it affects. \\
\hline
Export precondition & ONNX, AOT, quantization, serving, or accelerator-specific
constraints are violated. & Report the backend-specific contract and the nearest
portable replacement if known. \\
\hline
Unsupported dynamic region & The program constructs shapes or control flow from
values that the selected analysis mode cannot resolve. & Return \unknown{} with
the unsupported expression and a suggestion for an input hint or annotation. \\
\hline
Baseline-suppressed issue & A finding exists but is acknowledged by an adoption
baseline. & Keep the evidence in machine-readable output while not failing the
incremental gate. \\
\hline
Evidence inconsistency & A dashboard, paper artifact, corpus label, or coverage
table is stale relative to code or tests. & Fail regeneration checks and point
to the artifact that must be rebuilt. \\
\hline
\end{longtable}

\section{Reviewer Checklist}

A reviewer can inspect the trust boundary by reading
\texttt{SOUNDNESS\_CONTRACT.md}, \texttt{VERIFIABLE\_FRAGMENT.md}, and
\texttt{operator\_confidence\_table.json}.  They can run the latest padding
tests with:
\[
\texttt{python -m pytest -q tests/test\_pad\_semantics.py}.
\]
They can inspect the empirical story through \texttt{make reproduce}, the
dashboard tests, real benchmark loaders, baseline comparisons, and the
reproducibility directory.  They can inspect formal claims through the Lean
files and axiom-audit tests.  Most importantly, they can observe that unsupported
or symbolic-hard cases are not hidden: \tg{} either proves, refutes, or says
\unknown{} with a reason.

\end{document}
